% !TEX TS-program = lualatex

%%%
% src::
%   url = https://github.com/pfradin/luadraw/discussions/60
%%%

\documentclass[border = 3pt]{standalone}

%%%
% Des couleurs faciles d'emploi via le package \xcolor.
%%%
\usepackage[svgnames]{xcolor}

%%%
% La bibliothèque \luadraw allie une facilité d’utilisation
% à un rendu particulièrement soigné.
%%%
\usepackage[3d]{luadraw}

\begin{document}

\begin{luadraw*}{name = poly_diamond-in-sphere}
local ld = luadraw

local Origin = ld.pt3d.Origin
local M      = ld.pt3d.M

local vecI = ld.pt3d.vecI
local vecJ = ld.pt3d.vecJ
local vecK = ld.pt3d.vecK

------
-- prototype::
--     center   : XXXX
--     vertice  : XXXX
--     normal   : XXXX
--     nb_sides : XXXX
--
--     return: ???? un diamant dont les sommets sont sur une sphère.
------
local poly_diamond = function(
  center,
  vertice,
  normal,
  nb_sides
)
-- Construction du 1ier sommet sur l'axe de rotation.
  local r = ld.pt3d.abs(vertice - center)
  normal  = ld.pt3d.normalize(normal)
  local S = center + r*normal

-- ''2*center - S = center + (center - S)'' est le 2ième sommet
-- sur l'axe de rotation (symétrie de centre ''center'')
  local L = {S, vertice, 2*center - S}

-- Le diamant s'obtient en faisant tourner une ligne "latérale"
-- définie par les sommets sur l'axe de rotation et celui fourni
-- par l'utilisateur.
  local P = ld.rotline(
    L,                 -- Ligne "latérale" servant de base à la
                       -- construction du diamand.
    {center, normal},  -- Axe orienté de rotation.
    0, 360,            -- De `0` à `360` degrés.
    {nbdots = nb_sides + 1}
  )

  return ld.facet2poly(P)
end

------
-- ¨Def de la zone graphique.
------
local graphview = ld.graph3d:new{
  viewdir  = {30, 70},
  adjust2d = true,
  bbox     = false
}

------
-- Réglages liés aux lignes.
------
graphview:Linewidth(4)

ld.Hiddenlines     = true
ld.Hiddenlinestyle = "dashed"

------
-- Tracé de deux diamants.
------
graphview:Dpoly(
  poly_diamond(
    M(-4, -4, 0),
    M(-2, -2, 0),
    vecK,
    9
  ),
  {
-- ''ld.mShadedHidden'' demande la coloration des facettes selon
-- leur inclinaison et l'affichage des arêtes cachées.
    mode  = ld.mFlatHidden,
    color = "orange"
  }
)

graphview:Dpoly(
  poly_diamond(
    M(3, 3.5, 1), 
    M(3, 3, 3), 
    vecJ,
    7
  ),
  {

-- ''ld.mShadedHidden'' demande la coloration des facettes selon
-- leur inclinaison et sans l'affichage des arêtes cachées.
    mode  = ld.mFlat,
    color = "LightBlue"
  }
)

------
-- Montrons le résultat de notre oeuvre.
------
graphview:Show()
\end{luadraw*}

\end{document}
